\documentclass[../main.tex]{subfiles}

\begin{document}
	
	\chapter{Conclusioni}
	\label{ch:conclusioni}
	
	Il presente lavoro di tesi si è posto l'obiettivo di verificare in che misura i modelli \textit{data-driven} lineari possano offrire una reale interpretabilità fisica nella stima della viscosità dei liquidi sottoraffreddati. Assumendo come base la pipeline di estrazione delle proprietà chimico-fisiche proposta per il modello ViscNet~\cite{cassarViscNetNeuralNetwork2021}, si sono impiegati 560 descrittori per stimare i tre parametri dell'equazione MYEGA: la temperatura di transizione vetrosa ($T_g$), l'indice di fragilità ($m$) e la viscosità asintotica ad alta temperatura ($\log_{10}\eta_\infty$)~\cite{mauroViscosityGlassformingLiquids2009}.
	
	L'analisi parte dalla tesi secondo cui uno spazio descrittivo sufficientemente ricco consenta a modelli lineari di raggiungere un'accuratezza paragonabile a quella di architetture neurali complesse~\cite{sharmaInterpretabilityLinearRegression2026}. Tuttavia, in presenza di un numero elevato di variabili, la forte multicollinearità compromette la stabilità della regressione e ne oscura il significato fisico. Il confronto sistematico condotto sull'\textit{unseen test dataset} di 85 vetri ha permesso di delineare i limiti e le potenzialità delle diverse strategie adottate, misurando le performance sia sui singoli parametri termodinamici che sulla viscosità globale ricostruita:
	
	\begin{itemize}
		\item \textbf{Regressione Ridge:} Addestrata sull'intero spazio ad alta dimensionalità, conferma l'ipotesi della validità predittiva dei modelli lineari, ottenendo sulla viscosità globale ottime prestazioni ($R^2 = 0.920$, $\text{RMSE} = 1.147$) che competono direttamente con il benchmark della rete neurale ViscNet~\cite{cassarViscNetNeuralNetwork2021}. Ciononostante, come confermato dall'analisi del parametro di regolarizzazione $\alpha$, la penalizzazione $L_2$ non induce sparsità, spalmando il carico predittivo in una distribuzione continua di pesi che varia dinamicamente. Ciò impedisce di isolare le variabili fisicamente determinanti, riducendo di fatto la regressione a una "scatola nera" lineare.
		
		\item \textbf{Filtraggio VIF e OLS:} L'applicazione del \textit{Variance Inflation Factor} (VIF) ha abbattuto la collinearità riducendo lo spazio a 45 \textit{features}. Il modello ottiene prestazioni intermedie ($R^2 = 0.776$, $\text{RMSE} = 1.918$), ma si scontra con l'incapacità matematica dell'algoritmo OLS di forzare a zero le variabili meno rilevanti. La rimozione a priori della collinearità penalizza il potere predittivo senza restituire una gerarchia fisica ristretta e inequivocabile.
		
		\item \textbf{Regressione Elastic Net:} Grazie alla combinazione delle penalità $L_1$ ed $L_2$, Elastic Net supera l'arbitrarietà del VIF forzando a zero i coefficienti ridondanti. Pur registrando un netto calo predittivo globale ($R^2 = 0.670$), il modello estrae un nucleo coerente di descrittori fisici. Per la fragilità ($m$) dominano l'affinità elettronica e l'entalpia di fusione; per la $T_g$ emergono l'ingombro sterico, la valenza elettronica e i \textit{bandgap}. Al contrario, per $\log_{10}\eta_\infty$ le performance collassano ($R^2 = 0.098$): un risultato che non è un fallimento del modello, ma la conferma fisica della natura asintotica del parametro.
		
		\item \textbf{Principal Component Regression (PCR) e Supervised PCR (SPCR):} L'approccio basato sulla scomposizione in componenti principali risolve definitivamente la collinearità proiettando le \textit{features} su direzioni ortogonali. Se la PCR standard raggiunge un buon compromesso con 15 componenti ($R^2 = 0.762$), l'introduzione della Supervised PCR (SPCR) ottimizza l'allineamento con i target, ottenendo performance superiori sulla viscosità globale ($R^2 = 0.813$, $\text{RMSE} = 1.756$) utilizzando un sottospazio ancor più ristretto di sole 10 componenti. L'analisi di questi modelli evidenzia che le componenti responsabili della quota dominante della varianza geometrica dei dati (es. PC1 e PC2) risultano disaccoppiate dalla dinamica vetrosa. La reale capacità predittiva risiede in specifici modi collettivi a più bassa varianza~\cite{sharmaInterpretabilityLinearRegression2026}. Andando ad analizzare i carichi (\textit{loadings}) di queste componenti predittive (es. PC6 e PC11), si ritrovano proiettate le medesime proprietà chimico-fisiche evidenziate da Elastic Net (come l'entalpia di fusione per $m$, la valenza e l'ingombro sterico per $T_g$). Ciò fornisce una validazione incrociata e indipendente della reale rilevanza fisica di tali descrittori, confermando che l'effettiva dimensionalità del problema risiede in un numero limitato di gradi di libertà chimico-fisici.
	\end{itemize}

	In conclusione, questo studio dimostra che nei modelli \textit{data-driven} l'interpretabilità non è garantita automaticamente dalla linearità matematica, ma dipende criticamente dalla corretta gestione statistica della multicollinearità. Se l'obiettivo primario è la pura accuratezza, una regressione Ridge ad alta dimensionalità rappresenta un ottimo strumento eguagliando le prestazioni delle reti neurali. Quando invece lo scopo è l'estrazione di informazione scientifica, l'utilizzo di metodologie che riducono la dimensionalità e risolvono la collinearità (come SPCR ed Elastic Net) permette di individuare il compromesso ottimale: accettare un calo quantificabile nelle metriche di errore in cambio di una descrizione termodinamica più trasparente e robusta. Ulteriori indagini su specifici \textit{outliers} e sull'analisi dei residui (riportati nell'Appendice~\ref{app:grafici} per i vetri di test) potranno chiarire ulteriormente se i limiti predittivi riscontrati risiedano in carenze dello spazio descrittivo o in anomalie intrinseche dei dati sperimentali.
	
\end{document}